\uuid{RDeZ}
\titre{Intégration par parties/linéarité}
\auteur{hanani ; mzari}
\organisation{exo7}
\licence{CC BY-NC-SA 4.0}
\datecreate{2020-01-20}
\niveau{L1}
\module{Analyse}
\chapitre{Calcul d'intégrales}
\sousChapitre{Autre}
\difficulte{3}
\theme{Calculs d'intégrales, intégration par parties, linéarité}
\exo7qcmid{lille-2.integrales.031}

\contenu{
\begin{question}
\qclassification{127.99}
\qtags{motcle=intégration par parties/linéarité}
On pose $\displaystyle I=\int _0^{\pi}\mathrm{e}^x\cos ^2x\mathrm{d}x$, $\displaystyle J=\int _0^{\pi}\mathrm{e}^x\sin ^2x\mathrm{d}x$ et $\displaystyle K=\int _0^{\pi}\mathrm{e}^x\cos (2x)\mathrm{d}x$. Parmi les affirmations suivantes, cocher celles qui sont vraies :
\begin{answers}
\good{A l'aide de deux intégrations par parties successives, on obtient : $\displaystyle K=\mathrm{e}^{\pi}-1-4K$ et donc $\displaystyle K=\frac{\mathrm{e}^{\pi}-1}{5}$.}
\bad{$\displaystyle I+J=\mathrm{e}^{\pi}$ et $I-J=K$.}
\bad{$\displaystyle I=\frac{6\mathrm{e}^{\pi}-1}{10}$.}
\bad{$\displaystyle J=\frac{4\mathrm{e}^{\pi}+1}{10}$.}
\end{answers}
\vskip2mm
\begin{explanations}
Deux intégrations par parties successives, avec $u=\cos (2x)$ et $v=\mathrm{e}^x$ et ensuite avec $u=\sin (2x)$ et $v=\mathrm{e}^x$, donnent
$$K=\mathrm{e}^{\pi}-1+2\int _0^{\pi} \mathrm{e}^x\sin (2x)\mathrm{d}x=\mathrm{e}^{\pi}-1-4K\Rightarrow K=\frac{\mathrm{e}^{\pi}-1}{5}.$$
A l'aide des relations $\cos ^2x+\sin ^2x=1$ et  $\cos ^2x-\sin ^2x=\cos (2x)$, on obtient :
$$I+J=\mathrm{e}^{\pi}-1\mbox{ et }I-J=K\Rightarrow I=\frac{3(\mathrm{e}^{\pi}-1)}{5}\quad \mbox{et}\quad J=\frac{2(\mathrm{e}^{\pi}-1)}{5}.$$
\end{explanations}
\end{question}
}
